% =============================================================================
% APPENDIX DA2: REF-DEFENSE-MATRIX - Response Variants by Audience and Length
% =============================================================================
% LANGUAGE: german/english
% =============================================================================
%
% HEADER BLOCK
% -----------------------------------------------------------------------------
% Appendix:      DA2
% Category:      REF (Reference)
% Name:          DEFENSE-MATRIX
% Title:         Response Variants for Critical Questions
% Version:       1.0 (January 2026)
% Status:        ACTIVE
% Last Updated:  2026-01-30
% Authors:       EBF Development Team
% Dependencies:  DA (REF-DEFENSE), PR (CORE-PRINCIPLES), BBB (CORE-WHERE)
% Primary Chapter: 1 (Introduction)
% Secondary:     18 (Intervention Design)
% -----------------------------------------------------------------------------
%
% PURPOSE
% -----------------------------------------------------------------------------
% This appendix provides 6 response variants for each of the 20 critical
% questions documented in Appendix DA:
%   - 3 lengths: KURZ (1-2 sentences), MITTEL (paragraph), LANG (full)
%   - 2 audiences: WISSENSCHAFT (Ernst Fehr), PRAXIS (Middle Manager)
% =============================================================================

\chapter{Response Variants for Critical Questions}
\label{app:DA2}

%%%%%%%%%%%%%%%%%%%%%%%%%%%%%%%%%%%%%%%%%%%%%%%%%%%%%%%%%%%%%%%%%%%%%%%%%%%%%%%
% HEADER BLOCK
%%%%%%%%%%%%%%%%%%%%%%%%%%%%%%%%%%%%%%%%%%%%%%%%%%%%%%%%%%%%%%%%%%%%%%%%%%%%%%%

\begin{tcolorbox}[colback=gray!5!white, colframe=gray!75!black,
    title=Appendix Metadata]
\begin{tabular}{@{}ll@{}}
\textbf{Code:} & DA2 \\
\textbf{Category:} & REF (Reference) \\
\textbf{Name:} & DEFENSE-MATRIX \\
\textbf{Version:} & 1.0 (January 2026) \\
\textbf{Dependencies:} & DA, PR, BBB, AN, CCC/DDD \\
\textbf{Primary Chapter:} & Chapter 1 (Introduction) \\
\textbf{Secondary Chapters:} & Chapter 18 (Intervention Design), All \\
\end{tabular}
\end{tcolorbox}

%%%%%%%%%%%%%%%%%%%%%%%%%%%%%%%%%%%%%%%%%%%%%%%%%%%%%%%%%%%%%%%%%%%%%%%%%%%%%%%
% CHAPTER LINKAGE
%%%%%%%%%%%%%%%%%%%%%%%%%%%%%%%%%%%%%%%%%%%%%%%%%%%%%%%%%%%%%%%%%%%%%%%%%%%%%%%

\begin{tcolorbox}[colback=blue!5!white, colframe=blue!75!black,
    title=Chapter Linkage]
\textbf{Primary:} Chapter 1 (Introduction to EBF)---provides defense responses for framework presentation.

\textbf{Secondary:}
\begin{itemize}[nosep]
    \item Chapter 18: Intervention design context for practitioner questions
    \item All chapters: General framework defense applicable throughout
\end{itemize}

\textbf{Reading Path:} DA (Axioms) $\rightarrow$ DA2 (Response Variants) $\rightarrow$ PR (Principles)
\end{tcolorbox}

%%%%%%%%%%%%%%%%%%%%%%%%%%%%%%%%%%%%%%%%%%%%%%%%%%%%%%%%%%%%%%%%%%%%%%%%%%%%%%%
% SCOPE BOX
%%%%%%%%%%%%%%%%%%%%%%%%%%%%%%%%%%%%%%%%%%%%%%%%%%%%%%%%%%%%%%%%%%%%%%%%%%%%%%%

\begin{tcolorbox}[colback=green!5!white, colframe=green!75!black,
    title=Appendix Scope: Response Variants]
\textbf{Ziel:} Provide ready-to-use responses for defending EBF against critical questions.

\textbf{In-Scope:}
\begin{itemize}[nosep]
    \item 6 response variants per question (3 lengths $\times$ 2 audiences)
    \item Scientific citations for LANG versions
    \item Quick reference table for all 20 questions
\end{itemize}

\textbf{Out-of-Scope:}
\begin{itemize}[nosep]
    \item Formal axioms (see Appendix DA)
    \item Detailed parameter documentation (see Appendix BBB)
    \item LLMMC methodology details (see Appendix AN)
\end{itemize}

\textbf{Deliverables:}
\begin{itemize}[nosep]
    \item Response selection guide
    \item Audience profile table
    \item Full response variants for W1, P1, P6, W3
    \item Quick reference matrix for all 20 questions
\end{itemize}
\end{tcolorbox}

%%%%%%%%%%%%%%%%%%%%%%%%%%%%%%%%%%%%%%%%%%%%%%%%%%%%%%%%%%%%%%%%%%%%%%%%%%%%%%%
% FUNDAMENTAL QUESTION
%%%%%%%%%%%%%%%%%%%%%%%%%%%%%%%%%%%%%%%%%%%%%%%%%%%%%%%%%%%%%%%%%%%%%%%%%%%%%%%

\begin{tcolorbox}[colback=red!5!white, colframe=red!75!black,
    title=Fundamental Question]
\textbf{How do we defend EBF against critical questions from different audiences with appropriate depth and language?}

This appendix answers: Given a critical question (P1--P10 or W1--W10), what is the optimal response for a specific audience (scientific vs. practitioner) and context (elevator pitch vs. detailed discussion)?
\end{tcolorbox}

%%%%%%%%%%%%%%%%%%%%%%%%%%%%%%%%%%%%%%%%%%%%%%%%%%%%%%%%%%%%%%%%%%%%%%%%%%%%%%%
% ABSTRACT
%%%%%%%%%%%%%%%%%%%%%%%%%%%%%%%%%%%%%%%%%%%%%%%%%%%%%%%%%%%%%%%%%%%%%%%%%%%%%%%

\begin{abstract}
This appendix provides ready-to-use responses for the 20 critical questions
identified in Appendix DA. Each question has 6 response variants: 3 lengths
(short/medium/long) $\times$ 2 audiences (scientific/practitioner). Use the
appropriate variant based on context, time available, and audience expertise.
All scientific (LANG) versions include proper citations to the behavioral economics
literature (Kahneman, Tversky, Fehr, Thaler, Camerer, et al.).
\end{abstract}

%%%%%%%%%%%%%%%%%%%%%%%%%%%%%%%%%%%%%%%%%%%%%%%%%%%%%%%%%%%%%%%%%%%%%%%%%%%%%%%
\section{How to Use This Document}
\label{sec:da2-howto}
%%%%%%%%%%%%%%%%%%%%%%%%%%%%%%%%%%%%%%%%%%%%%%%%%%%%%%%%%%%%%%%%%%%%%%%%%%%%%%%

\begin{tcolorbox}[colback=blue!5!white, colframe=blue!75!black,
    title=Response Selection Guide]

\textbf{Step 1: Identify the Question}
\begin{itemize}[nosep]
    \item Practice questions: P1--P10 (utility, ROI, complexity)
    \item Science questions: W1--W10 (falsifiability, identification, validity)
\end{itemize}

\textbf{Step 2: Choose Audience}
\begin{itemize}[nosep]
    \item \textbf{WISSENSCHAFT} ($\mathcal{W}$): Peer reviewers, academics, Ernst Fehr
    \item \textbf{PRAXIS} ($\mathcal{P}$): Managers, consultants, clients
\end{itemize}

\textbf{Step 3: Choose Length}
\begin{itemize}[nosep]
    \item \textbf{KURZ} (K): 1--2 sentences, elevator pitch, immediate response
    \item \textbf{MITTEL} (M): 1 paragraph, meeting response, email
    \item \textbf{LANG} (L): Full explanation, presentation, written defense
\end{itemize}

\textbf{Notation:} Response variant = Question + Audience + Length

Example: W1-$\mathcal{W}$-M = Tautology question, scientific audience, medium length

\end{tcolorbox}

%%%%%%%%%%%%%%%%%%%%%%%%%%%%%%%%%%%%%%%%%%%%%%%%%%%%%%%%%%%%%%%%%%%%%%%%%%%%%%%
\section{Audience Profiles}
\label{sec:da2-audiences}
%%%%%%%%%%%%%%%%%%%%%%%%%%%%%%%%%%%%%%%%%%%%%%%%%%%%%%%%%%%%%%%%%%%%%%%%%%%%%%%

\begin{table}[htbp]
\centering
\caption{Target Audience Characteristics}
\renewcommand{\arraystretch}{1.4}
\begin{tabular}{@{}lp{6cm}p{6cm}@{}}
\toprule
& \textbf{WISSENSCHAFT ($\mathcal{W}$)} & \textbf{PRAXIS ($\mathcal{P}$)} \\
\midrule
\textbf{Archetyp} & Prof. Ernst Fehr & Abteilungsleiter, 45 Jahre \\
\textbf{Ausbildung} & PhD Economics, Top-5 Journals & MBA oder Berufserfahrung \\
\textbf{Methodik} & Kausalidentifikation, RCT, Peer Review & ``Was funktioniert in der Praxis'' \\
\textbf{Sprache} & Formal, präzise, Fachbegriffe & Konkret, Beispiele, Analogien \\
\textbf{Überzeugung durch} & Logik, Evidenz, Konsistenz & Erfolgsgeschichten, ROI, Einfachheit \\
\textbf{Skeptisch bei} & Unbegründeten Behauptungen & Theorie ohne Praxisbezug \\
\textbf{Zeitbudget} & Liest Paper vollständig & Liest Executive Summary \\
\bottomrule
\end{tabular}
\end{table}

%%%%%%%%%%%%%%%%%%%%%%%%%%%%%%%%%%%%%%%%%%%%%%%%%%%%%%%%%%%%%%%%%%%%%%%%%%%%%%%
%%%%%%%%%%%%%%%%%%%%%%%%%%%%%%%%%%%%%%%%%%%%%%%%%%%%%%%%%%%%%%%%%%%%%%%%%%%%%%%
\section{W1: Tautologie / ``If Everything Is Modelable...''}
\label{sec:da2-w1}
%%%%%%%%%%%%%%%%%%%%%%%%%%%%%%%%%%%%%%%%%%%%%%%%%%%%%%%%%%%%%%%%%%%%%%%%%%%%%%%
%%%%%%%%%%%%%%%%%%%%%%%%%%%%%%%%%%%%%%%%%%%%%%%%%%%%%%%%%%%%%%%%%%%%%%%%%%%%%%%

\begin{tcolorbox}[colback=red!5!white, colframe=red!75!black,
    title=W1: ``Wenn alles modellierbar ist, ist nichts falsifizierbar'']
\textbf{Kern:} Der fundamentalste wissenschaftliche Einwand---EBF sei eine Tautologie.
\end{tcolorbox}

%------------------------------------------------------------------------------
\subsection{W1: Wissenschaft ($\mathcal{W}$)}
%------------------------------------------------------------------------------

\begin{tcolorbox}[colback=green!5!white, colframe=green!75!black,
    title=W1-$\mathcal{W}$-K: Kurz (Wissenschaft)]

EBF ist falsifizierbar durch die Null-Hypothese $\gamma = 0$ (keine Komplementarität) und durch Vorhersagen mit Konfidenzintervallen, deren Verletzung das Modell widerlegt---analog zur Falsifizierbarkeit physikalischer Theorien durch quantitative Vorhersagen, nicht durch qualitative Flexibilität.

\end{tcolorbox}

\begin{tcolorbox}[colback=green!5!white, colframe=green!75!black,
    title=W1-$\mathcal{W}$-M: Mittel (Wissenschaft)]

Der Tautologie-Einwand verwechselt parametrische Flexibilität mit Unfalsifizierbarkeit. EBF ist falsifizierbar auf drei Ebenen: (1) Die Null-Hypothese $H_0: \gamma = 0$ ist stets testbar---wenn additive Modelle systematisch besser erklären als komplementäre, ist EBF widerlegt. (2) Parameter entstammen der Literatur mit epistemischen Tags (EMP/THR/LLM), nicht post-hoc Anpassung. (3) Jede Analyse produziert Vorhersagen mit CI; Outcomes ausserhalb des 95\%-Intervalls falsifizieren das spezifische Modell. Die Kritik träfe zu, wenn Parameter frei wählbar wären---sind sie aber nicht. Appendix S dokumentiert 10 explizit falsifizierbare Hypothesen mit spezifizierten Testbedingungen.

\end{tcolorbox}

\begin{tcolorbox}[colback=green!5!white, colframe=green!75!black,
    title=W1-$\mathcal{W}$-L: Lang (Wissenschaft)]

\textbf{Der Einwand formal:}

Die Kritik folgt der Popper'schen Falsifikationslogik \citep{popper1959logic}: Sei $\mathcal{M}$ eine Modellklasse mit Parameterraum $\Theta$. $\mathcal{M}$ ist tautologisch wenn:
\begin{equation}
\forall D \in \mathcal{D}: \exists \theta \in \Theta: P(D | \mathcal{M}, \theta) > \epsilon
\end{equation}
d.h., für jeden Datensatz existiert eine Parameterkonfiguration, die ihn ``erklärt''. Dies entspricht Lakatos' Kritik an ``degenerativen Forschungsprogrammen'' \citep{kuhn1962}.

\textbf{Warum EBF diese Bedingung nicht erfüllt:}

\textbf{1. Eingeschränkter Parameterraum:}
Parameter in EBF sind nicht frei wählbar, sondern durch externe Evidenz eingeschränkt:
\begin{equation}
\Theta_{\text{EBF}} = \{\theta : \theta \in \text{CI}_{95\%}(\hat{\theta}_{\text{Literatur}})\}
\end{equation}
Beispiel: $\lambda \in [1.5, 2.8]$ basierend auf der originalen Prospect Theory \citep{PAP-kahneman1979prospect,kahneman1991loss} und nachfolgenden Meta-Analysen. Ein Modell, das $\lambda = 5.0$ benötigt, ist falsifiziert, da ausserhalb des empirisch etablierten Intervalls.

\textbf{2. Testbare Null-Hypothese:}
\begin{equation}
H_0: \gamma_{ij} = 0 \quad \forall i,j \quad \text{(Faktoren unabhängig)}
\end{equation}
\begin{equation}
H_1: \exists i,j: \gamma_{ij} \neq 0 \quad \text{(EBF-Behauptung)}
\end{equation}
Wenn $H_0$ die Daten systematisch besser erklärt (via Likelihood-Ratio, BIC, Cross-Validation), ist EBF falsifiziert. Die Komplementaritäts-These ist direkt aus der empirischen Literatur zu Interaktionseffekten abgeleitet \citep{fehr2002altruistic,PAP-fischbacher2001conditional}.

\textbf{3. Prädiktive Validierung:}
Sei $\hat{Y}$ die EBF-Vorhersage mit $\text{CI}_{95\%}(\hat{Y}) = [\hat{Y}_L, \hat{Y}_U]$. Das Modell ist falsifiziert wenn:
\begin{equation}
Y_{\text{observed}} \notin [\hat{Y}_L, \hat{Y}_U]
\end{equation}
Dies erfordert Pre-Registration der Vorhersage vor Datenerhebung---analog zur Methodologie in \citet{camerer2018evaluating} für Replikationsstudien.

\textbf{4. Kreuzvalidierung über Kontexte:}
Parameter $\theta$, geschätzt in Kontext $\Psi_1$, müssen in Kontext $\Psi_2$ funktionieren (mit dokumentierter Kontext-Modulation). Systematisches Versagen der Übertragung falsifiziert die behauptete Universalität. Die Kontextabhängigkeit folgt der etablierten Literatur zu situativen Faktoren \citep{PAP-cialdini2006influence,loewenstein2003projection}.

\textbf{Explizite Falsifikationen (Appendix S):}
\begin{enumerate}[nosep]
    \item Cross-Level Spillover: $\Delta U^{L+1} = f(\Delta U^L)$ mit Lag $\tau \in [3,6]$ Monate
    \item Coherence Trap: $\gamma_{\text{cross}} < -0.3 \Rightarrow$ Entscheidungsblockade
    \item Identity Activation: Priming $IDN^L \Rightarrow \Delta \text{Cooperation} \in [15\%, 35\%]$ \citep{charness2012groups}
\end{enumerate}

\textbf{Epistemische Demut:}
Die Falsifizierbarkeit muss \textit{gelebt} werden. Wenn Anwender Parameter post-hoc anpassen oder CIs ignorieren, degeneriert EBF zur Tautologie in der Praxis. Deshalb ist P5 (Falsifiability) ein konstitutives Prinzip, nicht eine optionale Eigenschaft.

\textbf{Schlüsselreferenzen:}
\begin{itemize}[nosep]
    \item Prospect Theory: \citet{PAP-kahneman1979prospect}, \citet{kahneman1991loss}
    \item Falsifikationismus: \citet{popper1959logic}
    \item Soziale Präferenzen: \citet{fehr2002altruistic}, \citet{PAP-fischbacher2001conditional}
    \item Replikationsmethodologie: \citet{camerer2018evaluating}
\end{itemize}

\end{tcolorbox}

%------------------------------------------------------------------------------
\subsection{W1: Praxis ($\mathcal{P}$)}
%------------------------------------------------------------------------------

\begin{tcolorbox}[colback=yellow!5!white, colframe=yellow!75!black,
    title=W1-$\mathcal{P}$-K: Kurz (Praxis)]

EBF ist kein Bullshit-Generator, weil wir \textbf{vorher} sagen müssen, was passieren wird---mit Bandbreiten. Wenn wir falsch liegen, ist das dokumentiert. Berater ohne Framework können immer behaupten, sie hätten Recht---wir nicht.

\end{tcolorbox}

\begin{tcolorbox}[colback=yellow!5!white, colframe=yellow!75!black,
    title=W1-$\mathcal{P}$-M: Mittel (Praxis)]

Stellen Sie sich vor, ein Arzt sagt nach jeder Behandlung: ``Genau das habe ich erwartet.'' Das wäre verdächtig. EBF funktioniert anders: Wir müssen \textbf{vor} der Intervention sagen, was passieren wird---mit konkreten Zahlen und Bandbreiten (z.B. ``Teilnahme steigt um 25-35\%''). Wenn das Ergebnis 10\% oder 60\% ist, war unsere Vorhersage falsch, und das ist dokumentiert.

Der Unterschied zu ``normaler'' Beratung: Dort kann man immer nachträglich erklären, warum etwas funktioniert hat oder nicht. Bei EBF gibt es einen klaren Test: Lag die Realität in unserer vorhergesagten Bandbreite? Ja = Modell bestätigt. Nein = Modell anpassen. Das ist keine Ausrede-Maschine---das ist Wissenschaft.

\end{tcolorbox}

\begin{tcolorbox}[colback=yellow!5!white, colframe=yellow!75!black,
    title=W1-$\mathcal{P}$-L: Lang (Praxis)]

\textbf{Das Problem einfach erklärt:}

Kritiker sagen: ``EBF hat so viele Stellschrauben, dass ihr jedes Ergebnis erklären könnt. Das ist wie Horoskope---klingt immer passend, sagt aber nichts.''

Das wäre ein tödlicher Vorwurf, wenn er stimmen würde. Tut er aber nicht. Hier ist warum:

\textbf{1. Wir müssen uns vorher festlegen}

Bevor wir eine Intervention empfehlen, müssen wir sagen:
\begin{itemize}[nosep]
    \item ``Wir erwarten, dass Teilnahme von 40\% auf 55-65\% steigt''
    \item ``Der Effekt sollte nach 2-4 Wochen messbar sein''
    \item ``Die Wirkung ist bei Segment X stärker als bei Segment Y''
\end{itemize}

Das ist keine vage Aussage à la ``es wird besser''. Das sind konkrete Zahlen, die falsch sein können.

\textbf{2. Die Zahlen kommen nicht aus der Luft}

Wenn wir sagen ``Loss Aversion ist etwa 2.25'', dann kommt das aus Studien, die seit 40 Jahren repliziert werden. Wir können nicht einfach sagen ``Loss Aversion ist 10'', weil das unseren Daten nicht entspricht.

\textbf{Analogie Wettervorhersage:}

\begin{center}
\begin{tabular}{ll}
\textbf{Schlechte Vorhersage} & ``Morgen gibt es Wetter'' \\
\textbf{Gute Vorhersage} & ``Morgen 15-18°C, 30\% Regenwahrscheinlichkeit'' \\
\end{tabular}
\end{center}

Die erste ist immer ``richtig'' aber nutzlos. Die zweite kann falsch sein---und genau das macht sie wertvoll.

EBF macht die zweite Art von Vorhersagen.

\textbf{3. Wir dokumentieren unsere Fehler}

Im Intervention Registry steht nicht nur, was wir empfohlen haben, sondern auch, ob es funktioniert hat. Wenn wir systematisch daneben liegen, sieht man das. Bei traditioneller Beratung gibt es diese Transparenz nicht.

\textbf{4. Der praktische Test}

Fragen Sie Ihren Berater: ``Können Sie mir eine Vorhersage geben, bei der Sie sagen würden: Wenn X nicht eintritt, war meine Empfehlung falsch?''

Wenn die Antwort ``Nein'' oder Ausweichen ist---dann arbeiten Sie mit einer Ausrede-Maschine.

EBF kann diese Frage beantworten. Jedes Mal.

\textbf{Zusammenfassung:}

\begin{center}
\begin{tabular}{lll}
& \textbf{Horoskop/Tautologie} & \textbf{EBF} \\
\midrule
Vorhersage & Vage, nachträglich & Konkret, vorher \\
Kann falsch sein? & Nein & Ja (dokumentiert) \\
Zahlen woher? & Erfunden & Aus Forschung \\
Transparenz & Keine & Volle Dokumentation \\
\end{tabular}
\end{center}

\end{tcolorbox}

%%%%%%%%%%%%%%%%%%%%%%%%%%%%%%%%%%%%%%%%%%%%%%%%%%%%%%%%%%%%%%%%%%%%%%%%%%%%%%%
%%%%%%%%%%%%%%%%%%%%%%%%%%%%%%%%%%%%%%%%%%%%%%%%%%%%%%%%%%%%%%%%%%%%%%%%%%%%%%%
\section{P1: ``Das ist nur ein aufgemotzter Chatbot''}
\label{sec:da2-p1}
%%%%%%%%%%%%%%%%%%%%%%%%%%%%%%%%%%%%%%%%%%%%%%%%%%%%%%%%%%%%%%%%%%%%%%%%%%%%%%%
%%%%%%%%%%%%%%%%%%%%%%%%%%%%%%%%%%%%%%%%%%%%%%%%%%%%%%%%%%%%%%%%%%%%%%%%%%%%%%%

\begin{tcolorbox}[colback=red!5!white, colframe=red!75!black,
    title=P1: ``Am Ende fragt ihr einfach ChatGPT'']
\textbf{Kern:} Der häufigste Praxis-Einwand---EBF sei nur ein LLM mit Marketing.
\end{tcolorbox}

%------------------------------------------------------------------------------
\subsection{P1: Wissenschaft ($\mathcal{W}$)}
%------------------------------------------------------------------------------

\begin{tcolorbox}[colback=green!5!white, colframe=green!75!black,
    title=P1-$\mathcal{W}$-K: Kurz (Wissenschaft)]

EBF verwendet LLMs als Werkzeug innerhalb einer strukturierten Methodologie---analog zur Verwendung von Statistik-Software in ökonometrischer Analyse. Die Struktur (10C-Framework, Kontext-Spezifikation, Parameter-Restriktionen) geht der LLM-Nutzung voraus und determiniert sie.

\end{tcolorbox}

\begin{tcolorbox}[colback=green!5!white, colframe=green!75!black,
    title=P1-$\mathcal{W}$-M: Mittel (Wissenschaft)]

Der Einwand verwechselt Werkzeug mit Methodik. Ein Ökonometriker, der Stata verwendet, ist kein ``Stata-Operator''---die Methodologie (IV, DiD, RDD) determiniert, wie das Werkzeug eingesetzt wird. EBF verhält sich analog: Das 10C-Framework definiert die Analysestuktur, die $\Psi$-Spezifikation den Kontext, die Parameter-Datenbank (BBB) die zulässigen Werte. Das LLM wird dann verwendet für: (a) Prior-Generierung via LLMMC mit dokumentierter Unsicherheit, (b) Literatur-Synthese als Ausgangspunkt für manuelle Validierung, (c) Kommunikations-Generierung nach abgeschlossener Analyse. Der Unterschied zu ``einfach ChatGPT fragen'' ist strukturelle Präzedenz: Die Antwort wird durch das Framework determiniert, nicht durch das LLM generiert.

\end{tcolorbox}

\begin{tcolorbox}[colback=green!5!white, colframe=green!75!black,
    title=P1-$\mathcal{W}$-L: Lang (Wissenschaft)]

\textbf{Formale Unterscheidung:}

\textbf{Reines LLM:}
\begin{equation}
\text{Output} = \text{LLM}(\text{Prompt})
\end{equation}
Das Output ist eine Funktion des Prompts und des LLM-Trainings. Keine externe Struktur, keine Parameter-Restriktionen, keine Falsifizierbarkeit.

\textbf{EBF mit LLM-Unterstützung:}
\begin{equation}
\text{Output} = f\Big(\underbrace{\Psi}_{\text{Kontext}}, \underbrace{10C}_{\text{Struktur}}, \underbrace{\Theta_{\text{BBB}}}_{\text{Parameter}}, \underbrace{\text{LLM}(\cdot)}_{\text{Werkzeug}}\Big)
\end{equation}

\textbf{Wo LLM verwendet wird (mit Einschränkungen):}

\begin{enumerate}
    \item \textbf{LLMMC Prior-Generierung:}
    \begin{equation}
    \hat{\theta}_{\text{prior}} = \text{LLM}(\text{Context}, \text{Parameter-Definition})
    \end{equation}
    Mit expliziter Unsicherheitsquantifizierung und nachfolgendem Bayesian Updating gegen Literatur. Die Methodologie ist in Appendix AN (METHOD-LLMMC) formalisiert und folgt der Logik von Expert Elicitation \citep{ohaganelicitation}.

    \item \textbf{Literatur-Screening:}
    LLM identifiziert potenziell relevante Papers; finale Selektion und Extraktion manuell.

    \item \textbf{Output-Generierung:}
    Nach abgeschlossener Analyse: Formatierung für Zielgruppe (8D-Algorithmus, Appendix CCC).
\end{enumerate}

\textbf{Wo LLM NICHT verwendet wird:}

\begin{enumerate}
    \item \textbf{Struktur-Entscheidungen:} 10C-Dimensionen sind framework-definiert, nicht LLM-generiert. Die Utility-Dimensionen (FEPSDE) folgen \citet{fehr2002altruistic}.
    \item \textbf{Parameter-Werte:} Kommen aus BBB (Literatur), nicht aus LLM-``Wissen''. Beispiel: Loss Aversion $\lambda$ aus \citet{PAP-kahneman1979prospect}.
    \item \textbf{Kausal-Aussagen:} Identifikation erfordert Forschungsdesign \citep{PAP-camerer2003behavioral}, nicht LLM-Korrelationen.
    \item \textbf{Vorhersagen:} Emergieren aus Modell + Parameter, nicht aus LLM-Extrapolation.
\end{enumerate}

\textbf{Empirischer Test:}

Experiment: Dieselbe Frage an (a) ChatGPT direkt, (b) EBF-Workflow.

\textit{Frage:} ``Wie erhöhe ich die Sparquote bei 45-55-Jährigen in der Schweiz?''

\textit{ChatGPT:} Generische Liste (Automatisierung, Ziele setzen, Belohnungen) ohne Kontext-Spezifizität, Parameter oder Unsicherheit.

\textit{EBF:}
\begin{itemize}[nosep]
    \item Kontext: CH, $\Psi_I$ = Säule-3a-System, $\Psi_K$ = hohe Eigenverantwortung
    \item Modell: Default + Social Proof, $\gamma_{\text{default,social}} = +0.15$ \citep{johnson2003defaults,PAP-allcott2011social}
    \item Vorhersage: $\Delta\text{Participation} \in [22\%, 34\%]$ (95\% CI)
    \item Parameter: $\lambda_{\text{CH}} = 2.1$ (BFS-validiert), $\sigma_{\text{social}} = 0.12$ \citep{fehr2002cooperation}
\end{itemize}

Der Unterschied ist nicht graduell---er ist kategorial.

\textbf{Schlüsselreferenzen:}
\begin{itemize}[nosep]
    \item Behavioral Economics Framework: \citet{PAP-camerer2003behavioral}, \citet{thaler2015misbehaving}
    \item Default-Effekte: \citet{johnson2003defaults}, \citet{choi2004better}, \citet{benartzi2007save}
    \item Soziale Normen: \citet{PAP-cialdini2006influence}, \citet{PAP-allcott2011social}
    \item Kooperation: \citet{fehr2002cooperation}, \citet{PAP-fischbacher2001conditional}
\end{itemize}

\end{tcolorbox}

%------------------------------------------------------------------------------
\subsection{P1: Praxis ($\mathcal{P}$)}
%------------------------------------------------------------------------------

\begin{tcolorbox}[colback=yellow!5!white, colframe=yellow!75!black,
    title=P1-$\mathcal{P}$-K: Kurz (Praxis)]

ChatGPT gibt Antworten. EBF gibt Analysen. Der Unterschied: Bei ChatGPT wissen Sie nicht, woher die Antwort kommt. Bei EBF sehen Sie jeden Schritt, jede Zahl, jede Quelle---und können prüfen, ob es für Ihre Situation passt.

\end{tcolorbox}

\begin{tcolorbox}[colback=yellow!5!white, colframe=yellow!75!black,
    title=P1-$\mathcal{P}$-M: Mittel (Praxis)]

Machen Sie den Test: Fragen Sie ChatGPT ``Wie bringe ich meine Mitarbeiter zum Energiesparen?'' Sie bekommen eine nette Liste: ``1. Bewusstsein schaffen, 2. Anreize setzen, 3. Vorbildfunktion...'' Klingt gut, ist aber generisch.

Fragen Sie EBF dasselbe: Wir fragen zuerst zurück. Welches Land? Welche Branche? Welche Unternehmenskultur? Welche bisherigen Massnahmen? Dann bekommen Sie: ``In Ihrem Kontext (Schweizer Industrieunternehmen, hierarchische Kultur, bisherige Top-down-Kommunikation) empfehlen wir Social Proof + Defaults. Erwarteter Effekt: 15-25\% Reduktion. Quelle: 3 ähnliche Cases + Literatur.''

Das ist der Unterschied zwischen einer Antwort, die ``klingt gut'' und einer Analyse, die Sie verteidigen können.

\end{tcolorbox}

\begin{tcolorbox}[colback=yellow!5!white, colframe=yellow!75!black,
    title=P1-$\mathcal{P}$-L: Lang (Praxis)]

\textbf{Die Frage hinter der Frage:}

Wenn jemand sagt ``Das ist doch nur ChatGPT'', meint er eigentlich: ``Warum soll ich für etwas bezahlen, das ich kostenlos haben kann?''

Faire Frage. Hier ist die Antwort:

\textbf{Was ChatGPT macht:}
\begin{itemize}[nosep]
    \item Generiert Text basierend auf Trainingsmustern
    \item Keine Quellenangaben (oder halluzinierte)
    \item Keine Unsicherheits-Angaben (``könnte sein'' vs. ``ist wahrscheinlich'')
    \item Keine Kontextanpassung (Schweiz = USA = Nigeria)
    \item Keine Überprüfbarkeit (woher kommt diese Empfehlung?)
\end{itemize}

\textbf{Was EBF macht:}
\begin{itemize}[nosep]
    \item Strukturierte Analyse mit dokumentierten Schritten
    \item Jede Zahl hat eine Quelle (Paper, Studie, Datenbank)
    \item Explizite Unsicherheit (``25-35\%, nicht 30\%'')
    \item Kontextspezifisch (Schweiz $\neq$ USA $\neq$ Nigeria)
    \item Vollständig überprüfbar (Sie können jeden Schritt nachvollziehen)
\end{itemize}

\textbf{Analogie: Google Maps vs. Ortskundiger Taxifahrer}

\begin{center}
\begin{tabular}{lll}
& \textbf{Google Maps} & \textbf{Ortskundiger Taxifahrer} \\
\midrule
Schnell? & Ja & Ja \\
Günstig? & Ja & Nein \\
Immer verfügbar? & Ja & Nein \\
Kennt Baustellen heute? & Vielleicht & Ja \\
Kennt den Hintereingang? & Nein & Ja \\
Kann Kontext nutzen? & Begrenzt & Ja \\
\end{tabular}
\end{center}

ChatGPT ist Google Maps: Schnell, billig, oft gut genug.

EBF ist der ortskundige Fahrer: Teurer, aber kennt die Abkürzungen, die Baustellen, und warum Sie heute besser nicht über die Hauptstrasse fahren sollten.

\textbf{Wann brauchen Sie was?}

\begin{itemize}[nosep]
    \item \textbf{Einfache Frage, niedriges Risiko:} ChatGPT reicht
    \item \textbf{Wichtige Entscheidung, hohe Kosten bei Fehler:} EBF
    \item \textbf{Müssen Sie die Entscheidung verteidigen?} EBF
    \item \textbf{Einmalig vs. wiederholbar:} ChatGPT vs. EBF
\end{itemize}

\textbf{Der entscheidende Test:}

Können Sie mit der ChatGPT-Antwort in ein Board-Meeting gehen und sagen: ``Hier ist meine Empfehlung, hier sind die Quellen, hier ist die erwartete Wirkung, hier ist das Risiko wenn es nicht funktioniert''?

Wenn nein---dann brauchen Sie mehr als ChatGPT.

\textbf{Zusammenfassung:}

\begin{center}
\begin{tabular}{lcc}
& \textbf{ChatGPT} & \textbf{EBF} \\
\midrule
Antwort & $\checkmark$ & $\checkmark$ \\
Analyse & $\times$ & $\checkmark$ \\
Quellen & $\times$ & $\checkmark$ \\
Unsicherheit & $\times$ & $\checkmark$ \\
Kontext & $\times$ & $\checkmark$ \\
Verteidigbar & $\times$ & $\checkmark$ \\
\end{tabular}
\end{center}

\end{tcolorbox}

%%%%%%%%%%%%%%%%%%%%%%%%%%%%%%%%%%%%%%%%%%%%%%%%%%%%%%%%%%%%%%%%%%%%%%%%%%%%%%%
%%%%%%%%%%%%%%%%%%%%%%%%%%%%%%%%%%%%%%%%%%%%%%%%%%%%%%%%%%%%%%%%%%%%%%%%%%%%%%%
\section{P6: ``Die Parameter wirken willkürlich''}
\label{sec:da2-p6}
%%%%%%%%%%%%%%%%%%%%%%%%%%%%%%%%%%%%%%%%%%%%%%%%%%%%%%%%%%%%%%%%%%%%%%%%%%%%%%%
%%%%%%%%%%%%%%%%%%%%%%%%%%%%%%%%%%%%%%%%%%%%%%%%%%%%%%%%%%%%%%%%%%%%%%%%%%%%%%%

\begin{tcolorbox}[colback=red!5!white, colframe=red!75!black,
    title=P6: ``$\lambda = 2.25$? Woher kommt das?'']
\textbf{Kern:} Zweifel an der empirischen Fundierung der Parameter.
\end{tcolorbox}

%------------------------------------------------------------------------------
\subsection{P6: Wissenschaft ($\mathcal{W}$)}
%------------------------------------------------------------------------------

\begin{tcolorbox}[colback=green!5!white, colframe=green!75!black,
    title=P6-$\mathcal{W}$-K: Kurz (Wissenschaft)]

Jeder Parameter trägt einen epistemischen Tag (EMP/THR/LLM/ILL/HYP), der Herkunft und Unsicherheit dokumentiert. $\lambda = 2.25$ ist EMP (Tversky \& Kahneman 1992), mit Meta-Analyse-CI $[1.5, 2.8]$. Sensitivitätsanalyse quantifiziert die Ergebnisabhängigkeit von Parameter-Unsicherheit.

\end{tcolorbox}

\begin{tcolorbox}[colback=green!5!white, colframe=green!75!black,
    title=P6-$\mathcal{W}$-M: Mittel (Wissenschaft)]

Die Parameter-Kritik ist methodologisch berechtigt und wird durch ein dreistufiges Transparenz-System adressiert:

\textbf{1. Epistemische Tags:}
\begin{itemize}[nosep]
    \item EMP: Empirisch geschätzt (Studie zitiert, SE angegeben)
    \item THR: Theoretisch abgeleitet (Modell referenziert)
    \item LLM: LLMMC-Prior (Unsicherheit explizit, Kalibrierung dokumentiert)
    \item ILL: Illustrativ (nur für Beispiele, nicht für Schlussfolgerungen)
    \item HYP: Hypothetisch (Szenario-Analyse)
\end{itemize}

\textbf{2. Literatur-Verankerung:}
$\lambda = 2.25$ basiert auf \citet{PAP-kahneman1979prospect}, repliziert in >100 Studien. Meta-Analyse-Schätzung: $\hat{\lambda} = 2.25$ mit 95\%-CI $[1.93, 2.57]$.

\textbf{3. Sensitivitätsanalyse:}
Für jeden kritischen Parameter: $\partial \text{Result} / \partial \theta$. Wenn Ergebnis robust über CI---Parameter unkritisch. Wenn sensitiv---explizit als Unsicherheitsquelle dokumentiert.

\end{tcolorbox}

\begin{tcolorbox}[colback=green!5!white, colframe=green!75!black,
    title=P6-$\mathcal{W}$-L: Lang (Wissenschaft)]

\textbf{Das epistemische Problem:}

Die Kritik ``Parameter wirken willkürlich'' reflektiert ein legitimes Bedenken: Wenn Parameter frei wählbar sind, können Modelle jeden Datensatz post-hoc erklären, was Falsifizierbarkeit untergräbt \citep{popper1959logic}.

\textbf{EBF-Lösung: Vierstufiges Transparenz-System}

\textbf{1. Epistemische Tags (jeder Parameter trägt Herkunftsnachweis):}

\begin{center}
\begin{tabular}{lll}
\textbf{Tag} & \textbf{Bedeutung} & \textbf{Vertrauensniveau} \\
\midrule
EMP & Empirisch geschätzt, Studie zitiert, SE angegeben & Hoch \\
THR & Theoretisch abgeleitet, Modell referenziert & Mittel-Hoch \\
LLM & LLMMC-Prior, Unsicherheit explizit quantifiziert & Mittel \\
ILL & Illustrativ, nur für Beispiele & Niedrig \\
HYP & Hypothetisch, Szenario-Analyse & N/A \\
\end{tabular}
\end{center}

\textbf{2. Literatur-Verankerung für Kernparameter:}

\textit{Loss Aversion ($\lambda$):}
\begin{itemize}[nosep]
    \item Originalschätzung: $\lambda \approx 2.25$ \citep{PAP-kahneman1979prospect}
    \item Replikationen: >100 Studien in verschiedenen Domänen \citep{kahneman1991loss}
    \item Kontextvariation: $\lambda \in [1.5, 2.8]$ je nach Stakes und Framing
    \item Verwendung: Immer mit CI, nie als Punktschätzung
\end{itemize}

\textit{Discount Rate ($\beta$):}
\begin{itemize}[nosep]
    \item Quasi-hyperbolisch: $\beta \in [0.7, 0.9]$ \citep{laibson1997golden,PAP-odonoghue1999doing}
    \item Kontextabhängig: Gesundheit vs. Finanzen zeigen unterschiedliche $\beta$
\end{itemize}

\textit{Social Preference Parameter ($\sigma$):}
\begin{itemize}[nosep]
    \item Inequity Aversion: $\alpha \approx 0.85$, $\beta \approx 0.30$ \citep{fehr1993fairness}
    \item Conditional Cooperation: $\sigma \in [0.10, 0.25]$ \citep{PAP-fischbacher2001conditional}
\end{itemize}

\textbf{3. Sensitivitätsanalyse (obligatorisch):}

Für jeden kritischen Parameter $\theta$:
\begin{equation}
\text{Sensitivität} = \frac{\partial \hat{Y}}{\partial \theta} \cdot \frac{\sigma_\theta}{\hat{Y}}
\end{equation}

Klassifikation:
\begin{itemize}[nosep]
    \item Sensitivität $< 0.1$: Robust---Parameter-Unsicherheit unkritisch
    \item Sensitivität $\in [0.1, 0.3]$: Moderat---Dokumentieren
    \item Sensitivität $> 0.3$: Kritisch---Prominente Warnung
\end{itemize}

\textbf{4. Bayesian Updating gegen neue Evidenz:}

\begin{equation}
P(\theta | D_{\text{neu}}) \propto P(D_{\text{neu}} | \theta) \cdot P(\theta | D_{\text{Literatur}})
\end{equation}

Neue Studien aktualisieren bestehende Schätzungen, nicht ersetzen sie ad-hoc.

\textbf{Kontrast zu ``willkürlich'':}

\begin{center}
\begin{tabular}{lll}
& \textbf{Willkürlich} & \textbf{EBF} \\
\midrule
Herkunft & Unklar & Tagged (EMP/THR/LLM) \\
Wert & Punktschätzung & Intervall mit CI \\
Sensitivität & Ignoriert & Quantifiziert \\
Neue Evidenz & Ignoriert & Bayesian Update \\
Dokumentation & Keine & Vollständig (BBB) \\
\end{tabular}
\end{center}

\textbf{Schlüsselreferenzen:}
\begin{itemize}[nosep]
    \item Prospect Theory: \citet{PAP-kahneman1979prospect}, \citet{kahneman1991loss}
    \item Hyperbolisches Diskontieren: \citet{laibson1997golden}, \citet{PAP-odonoghue1999doing}
    \item Soziale Präferenzen: \citet{fehr1993fairness}, \citet{PAP-fischbacher2001conditional}, \citet{fehr2002altruistic}
    \item Methodik: \citet{PAP-camerer2003behavioral}, \citet{camerer2018evaluating}
\end{itemize}

\end{tcolorbox}

%------------------------------------------------------------------------------
\subsection{P6: Praxis ($\mathcal{P}$)}
%------------------------------------------------------------------------------

\begin{tcolorbox}[colback=yellow!5!white, colframe=yellow!75!black,
    title=P6-$\mathcal{P}$-K: Kurz (Praxis)]

Jede Zahl hat einen Stempel: ``Aus Studie'' (verlässlich), ``Aus Theorie'' (begründet), oder ``Geschätzt'' (mit Vorsicht). Sie können immer fragen: ``Woher kommt diese Zahl?''---und wir können es Ihnen zeigen.

\end{tcolorbox}

\begin{tcolorbox}[colback=yellow!5!white, colframe=yellow!75!black,
    title=P6-$\mathcal{P}$-M: Mittel (Praxis)]

Stellen Sie sich vor, Ihr Arzt sagt: ``Nehmen Sie 500mg.'' Sie fragen: ``Warum 500?'' Gute Antwort: ``Das ist die Standard-Dosierung aus klinischen Studien für Ihr Körpergewicht.'' Schlechte Antwort: ``Das machen wir immer so.''

EBF funktioniert wie der gute Arzt: Wenn wir sagen ``Loss Aversion ist 2.25'', dann heisst das: ``Menschen empfinden Verluste etwa 2,25-mal so stark wie gleich grosse Gewinne. Das kommt aus Studien mit tausenden Teilnehmern seit 1979, repliziert in über 100 Experimenten weltweit.''

Und wenn die Zahl für Ihre spezifische Situation anders sein könnte? Dann sagen wir das auch: ``Für Ihre Zielgruppe könnte es zwischen 1.8 und 2.8 liegen---hier ist, was das für die Empfehlung bedeutet.''

\end{tcolorbox}

%%%%%%%%%%%%%%%%%%%%%%%%%%%%%%%%%%%%%%%%%%%%%%%%%%%%%%%%%%%%%%%%%%%%%%%%%%%%%%%
%%%%%%%%%%%%%%%%%%%%%%%%%%%%%%%%%%%%%%%%%%%%%%%%%%%%%%%%%%%%%%%%%%%%%%%%%%%%%%%
\section{W3: ``Kein RCT für das Framework''}
\label{sec:da2-w3}
%%%%%%%%%%%%%%%%%%%%%%%%%%%%%%%%%%%%%%%%%%%%%%%%%%%%%%%%%%%%%%%%%%%%%%%%%%%%%%%
%%%%%%%%%%%%%%%%%%%%%%%%%%%%%%%%%%%%%%%%%%%%%%%%%%%%%%%%%%%%%%%%%%%%%%%%%%%%%%%

\begin{tcolorbox}[colback=red!5!white, colframe=red!75!black,
    title=W3: ``Wo ist der RCT, der zeigt, dass EBF besser ist?'']
\textbf{Kern:} Fehlende Framework-Level-Evidenz.
\end{tcolorbox}

%------------------------------------------------------------------------------
\subsection{W3: Wissenschaft ($\mathcal{W}$)}
%------------------------------------------------------------------------------

\begin{tcolorbox}[colback=green!5!white, colframe=green!75!black,
    title=W3-$\mathcal{W}$-K: Kurz (Wissenschaft)]

Der Einwand ist methodologisch berechtigt: Framework-Level-RCTs sind nicht verfügbar. Die Evidenzbasis operiert auf Komponenten-Ebene (2,328 integrierte Studien) und Case-Ebene (852 dokumentierte Anwendungen). Ein Cluster-RCT ``EBF-Beratung vs. Kontrolle'' wäre wünschenswert, ist aber methodologisch und praktisch herausfordernd.

\end{tcolorbox}

\begin{tcolorbox}[colback=green!5!white, colframe=green!75!black,
    title=W3-$\mathcal{W}$-M: Mittel (Wissenschaft)]

Wir unterscheiden drei Evidenz-Ebenen:

\textbf{Tier 1 (Stark):} Die 2,328 integrierten Papers liefern RCT-Evidenz für einzelne Mechanismen (Loss Aversion \citep{PAP-kahneman1979prospect}, Social Proof \citep{PAP-cialdini2006influence,PAP-allcott2011social}, Defaults \citep{johnson2003defaults,choi2004better}). Diese Evidenz ist robust.

\textbf{Tier 2 (Moderat):} 852 Cases dokumentieren Anwendungen mit Outcomes. Selection Bias ist wahrscheinlich (Erfolge überrepräsentiert), aber die Dokumentation ermöglicht Aggregation.

\textbf{Tier 3 (Fehlend):} Direkter RCT ``Beratung mit EBF'' vs. ``Beratung ohne EBF'' existiert nicht.

\textbf{Methodologische Herausforderungen für Tier-3-Evidenz:}
\begin{itemize}[nosep]
    \item Randomisierungseinheit: Berater? Projekt? Unternehmen?
    \item Verblindung: Nicht möglich (Berater weiss, ob er EBF nutzt)
    \item Outcome-Messung: Welcher Erfolgsindikator? Wann gemessen?
    \item Kontamination: Wissen diffundiert zwischen Gruppen
\end{itemize}

\textbf{Geplante Evaluation:} Pre-Registration von Vorhersagen + systematisches Outcome-Tracking als pragmatische Alternative zu RCT, analog zur Methodologie in \citet{camerer2018evaluating}.

\end{tcolorbox}

\begin{tcolorbox}[colback=green!5!white, colframe=green!75!black,
    title=W3-$\mathcal{W}$-L: Lang (Wissenschaft)]

\textbf{Der Einwand präzisiert:}

Die Forderung nach einem Framework-Level-RCT ist methodologisch korrekt. Um zu zeigen, dass ``EBF-basierte Beratung'' besser ist als ``Beratung ohne EBF'', wäre ein Experiment erforderlich wie:

\begin{equation}
\text{RCT}: \quad Y_{\text{EBF}} - Y_{\text{Control}} = \tau_{\text{Framework}} \quad \text{(zu schätzen)}
\end{equation}

Dieses Experiment existiert nicht. Wir sind transparent darüber.

\textbf{Die Evidenz-Hierarchie:}

\textbf{Tier 1: Komponenten-Evidenz (stark)}

Für einzelne Mechanismen existiert robuste RCT-Evidenz:

\begin{center}
\begin{tabular}{lll}
\textbf{Mechanismus} & \textbf{Effektgrösse} & \textbf{Evidenzquelle} \\
\midrule
Defaults & $+25\%$ bis $+35\%$ & \citet{johnson2003defaults}, \citet{choi2004better} \\
Social Proof & $+2\%$ bis $+10\%$ & \citet{PAP-allcott2011social}, \citet{PAP-schultz2007normative} \\
Loss Framing & $+5\%$ bis $+15\%$ & \citet{PAP-kahneman1979prospect}, \citet{benartzi1995myopic} \\
Commitment & $+10\%$ bis $+20\%$ & \citet{benartzi2007save}, \citet{bryan2010commitment} \\
\end{tabular}
\end{center}

Diese Effekte sind in hunderten von RCTs repliziert und robust über Kontexte \citep{PAP-camerer2003behavioral,thaler2015misbehaving}.

\textbf{Tier 2: Case-Evidenz (moderat)}

852 dokumentierte Anwendungsfälle im EBF Case Registry. Limitation: Selection Bias wahrscheinlich (Erfolge überrepräsentiert). Mitigation: Systematische Dokumentation aller Projekte, nicht nur Erfolge.

\textbf{Tier 3: Framework-Evidenz (fehlend)}

Kein RCT auf Ebene ``EBF vs. Nicht-EBF''.

\textbf{Warum Framework-RCTs schwierig sind:}

\begin{enumerate}
    \item \textbf{Randomisierungseinheit:} Berater? Projekt? Unternehmen? Jede Wahl hat Konsequenzen für Power und Interpretation.

    \item \textbf{Verblindung unmöglich:} Berater weiss, ob er EBF nutzt. Placebo-Effekte nicht ausschliessbar.

    \item \textbf{Outcome-Definition:} Was ist ``Erfolg''? Kundenzufriedenheit? Verhaltensänderung? ROI? Verschiedene Metriken können divergieren.

    \item \textbf{Zeithorizont:} Wann messen? Kurzfristige Effekte können sich von langfristigen unterscheiden \citep{benartzi2017should}.

    \item \textbf{Kontamination:} In Beratungsmarkt diffundiert Wissen. Kontrolle kann EBF-Elemente aufnehmen.
\end{enumerate}

\textbf{Analogie: Evidenzbasierte Medizin}

Die Frage ``Wirkt evidenzbasierte Medizin?'' wurde auch nicht durch einen RCT beantwortet. Stattdessen:
\begin{itemize}[nosep]
    \item Komponenten sind durch RCTs validiert (einzelne Therapien)
    \item Framework organisiert die Evidenz
    \item Pragmatische Evaluation durch Outcome-Tracking
\end{itemize}

EBF folgt demselben Ansatz.

\textbf{Pragmatische Evaluation (implementiert):}

\begin{enumerate}
    \item \textbf{Pre-Registration:} Jede Vorhersage wird vor Intervention dokumentiert
    \item \textbf{Outcome-Tracking:} Systematische Erfassung, ob Vorhersage eintraf
    \item \textbf{Deviation Analysis:} Bei Abweichung: Warum? Modell-Update?
    \item \textbf{Transparenz:} Case Registry enthält Erfolge UND Misserfolge
\end{enumerate}

Dies ist keine perfekte kausale Identifikation, aber besser als die Black-Box traditioneller Beratung.

\textbf{Offene Forderung:}

Wir fordern andere Beratungsansätze (McKinsey, BCG, traditionelle Praxis) auf, dieselbe Transparenz zu zeigen. Die asymmetrische Evidenzforderung---RCT für EBF, aber nicht für Alternativen---ist wissenschaftlich nicht gerechtfertigt.

\textbf{Schlüsselreferenzen:}
\begin{itemize}[nosep]
    \item Komponenten-Evidenz: \citet{PAP-kahneman1979prospect}, \citet{johnson2003defaults}, \citet{PAP-allcott2011social}
    \item Defaults in Retirement: \citet{choi2004better}, \citet{benartzi2007save}
    \item Methodologie: \citet{PAP-camerer2003behavioral}, \citet{camerer2018evaluating}
    \item Überblick: \citet{thaler2015misbehaving}, \citet{sunstein2015simplifying}
\end{itemize}

\end{tcolorbox}

%------------------------------------------------------------------------------
\subsection{W3: Praxis ($\mathcal{P}$)}
%------------------------------------------------------------------------------

\begin{tcolorbox}[colback=yellow!5!white, colframe=yellow!75!black,
    title=W3-$\mathcal{P}$-K: Kurz (Praxis)]

Haben wir einen wissenschaftlichen Beweis, dass EBF-Beratung besser ist als andere? Noch nicht---EBF ist 2 Jahre alt. Aber: Die Bausteine (Loss Aversion, Defaults, Social Proof) sind hundertfach bewiesen. Wir bauen auf solider Forschung auf.

\end{tcolorbox}

\begin{tcolorbox}[colback=yellow!5!white, colframe=yellow!75!black,
    title=W3-$\mathcal{P}$-M: Mittel (Praxis)]

Ehrliche Antwort: Wir haben keinen wissenschaftlichen Beweis, dass ``Beratung mit EBF'' besser ist als ``Beratung ohne EBF''. So ein Beweis würde erfordern, dass wir 100 Unternehmen zufällig in zwei Gruppen teilen und Jahre später vergleichen. Das hat niemand gemacht---nicht für EBF, nicht für McKinsey, nicht für BCG.

Was wir haben:
\begin{itemize}[nosep]
    \item 2,328 Studien, die zeigen, dass die Bausteine funktionieren
    \item 852 dokumentierte Projekte mit Ergebnissen
    \item Ein System, das Vorhersagen macht und dokumentiert, ob sie eintreffen
\end{itemize}

Fragen Sie mal McKinsey: ``Haben Sie einen RCT, der zeigt, dass Ihre Beratung wirkt?'' Die Antwort wird Schweigen sein.

Wir sind wenigstens transparent darüber, was wir wissen und was nicht.

\end{tcolorbox}

%%%%%%%%%%%%%%%%%%%%%%%%%%%%%%%%%%%%%%%%%%%%%%%%%%%%%%%%%%%%%%%%%%%%%%%%%%%%%%%
\section{Quick Reference: All 20 Questions}
\label{sec:da2-quickref}
%%%%%%%%%%%%%%%%%%%%%%%%%%%%%%%%%%%%%%%%%%%%%%%%%%%%%%%%%%%%%%%%%%%%%%%%%%%%%%%

\begin{table}[htbp]
\centering
\caption{Response Selection Matrix}
\renewcommand{\arraystretch}{1.2}
\scriptsize
\begin{tabular}{@{}clp{4cm}p{4cm}@{}}
\toprule
\textbf{\#} & \textbf{Frage} & \textbf{$\mathcal{W}$-K (Wissenschaft kurz)} & \textbf{$\mathcal{P}$-K (Praxis kurz)} \\
\midrule
\multicolumn{4}{@{}l}{\textit{\textbf{PRAXIS (P1--P10)}}} \\
P1 & Nur Chatbot & Strukturelle Präzedenz, LLM als Werkzeug & ChatGPT gibt Antworten, EBF gibt Analysen \\
P2 & Wo ist ROI? & Komponenten + Case + Framework Evidence & Die Bausteine sind bewiesen; Framework-Beweis kommt \\
P3 & Branche anders & Universelle Struktur, lokale $\Psi$ & 10C gilt überall; Ihr Kontext ist spezifisch \\
P4 & Zu komplex & Drei Modi: SCHNELL/STANDARD/TIEF & Komplexität ist optional, nicht Pflicht \\
P5 & Berater machten's & P7 Cumulation: Wissen akkumuliert & Bei uns bleibt das Wissen, bei denen geht es \\
P6 & Parameter willkürlich & Epistemische Tags, Sensitivität & Jede Zahl hat einen Stempel und eine Quelle \\
P7 & KI macht's schon & Analyse vs. Antwort, Struktur & ChatGPT = Antwort, EBF = verteidigbare Analyse \\
P8 & Nur Theorie & 852 reale Cases dokumentiert & Das sind echte Projekte, keine Gedankenspiele \\
P9 & Zu teuer & Total Cost of Intuition & Was kostet ein falscher Nudge bei 100k Kunden? \\
P10 & Kunde versteht's nicht & Analyse $\neq$ Kommunikation & Komplexe Analyse, einfache Präsentation \\
\addlinespace
\multicolumn{4}{@{}l}{\textit{\textbf{WISSENSCHAFT (W1--W10)}}} \\
W1 & Tautologie & $H_0: \gamma=0$ testbar, CI-Vorhersagen & Wir müssen vorher sagen, was passiert \\
W2 & $\gamma$ nicht identifiziert & Deskriptiv vs. kausal transparent & Wir sagen, was sicher ist und was nicht \\
W3 & Kein Framework-RCT & Tier-1/2/3 Evidenz-Hierarchie & Die Bausteine sind bewiesen; Rest kommt \\
W4 & Nur Meta-Framework & Framework ermöglicht neue Vorhersagen & Wie Periodensystem: organisiert + ermöglicht \\
W5 & LLMMC ungültig & Prior, nicht Posterior; Bayesian Update & Startpunkt, nicht Endpunkt; Literatur dominiert \\
W6 & Kein Peer Review & Multi-Channel; Papers in Arbeit & 2,328 Papers sind peer-reviewed \\
W7 & Normativ/Positiv vermischt & Strukturelle Trennung Ch.1-16 vs. 17-22 & Analyse getrennt von Empfehlung \\
W8 & Axiome nicht minimal & Operational vs. mathematisch minimal & Klarheit für Anwender $>$ Eleganz für Theoretiker \\
W9 & $\Psi$ nicht messbar & BCM2: 404 Faktoren operationalisiert & Wir haben es gemessen: 404 Faktoren für CH \\
W10 & Selection Bias & Acknowledged; Mitigation in progress & Wir dokumentieren auch Misserfolge \\
\bottomrule
\end{tabular}
\end{table}

%%%%%%%%%%%%%%%%%%%%%%%%%%%%%%%%%%%%%%%%%%%%%%%%%%%%%%%%%%%%%%%%%%%%%%%%%%%%%%%
\section{Cross-References}
\label{sec:da2-crossref}
%%%%%%%%%%%%%%%%%%%%%%%%%%%%%%%%%%%%%%%%%%%%%%%%%%%%%%%%%%%%%%%%%%%%%%%%%%%%%%%

\begin{itemize}[nosep]
    \item \textbf{DA (REF-DEFENSE):} Vollständige Axiome DA-1 bis DA-20
    \item \textbf{PR (CORE-PRINCIPLES):} Die 7 Prinzipien (P1-P7), die verteidigt werden
    \item \textbf{S (PREDICT-MASTER):} Die 10 falsifizierbaren Hypothesen
    \item \textbf{BBB (CORE-WHERE):} Parameter-Datenbank mit epistemischen Tags
    \item \textbf{AN (METHOD-LLMMC):} LLM Monte Carlo Methodologie
    \item \textbf{CCC/DDD (DOCTYPE):} 8D-Algorithmus für zielgruppengerechte Kommunikation
\end{itemize}

%%%%%%%%%%%%%%%%%%%%%%%%%%%%%%%%%%%%%%%%%%%%%%%%%%%%%%%%%%%%%%%%%%%%%%%%%%%%%%%
\section{Summary}
\label{sec:da2-summary}
%%%%%%%%%%%%%%%%%%%%%%%%%%%%%%%%%%%%%%%%%%%%%%%%%%%%%%%%%%%%%%%%%%%%%%%%%%%%%%%

\begin{tcolorbox}[colback=blue!5!white, colframe=blue!75!black,
    title=Key Takeaways]
\begin{enumerate}
    \item \textbf{Response Selection:} Choose variant based on audience ($\mathcal{W}$ vs. $\mathcal{P}$) and context (K/M/L).

    \item \textbf{Scientific Rigor:} LANG versions for $\mathcal{W}$ include proper citations to behavioral economics literature.

    \item \textbf{Practical Utility:} KURZ versions for $\mathcal{P}$ provide immediate, defensible responses.

    \item \textbf{Consistency:} All responses derive from the same underlying axioms (DA-0 to DA-20).

    \item \textbf{Transparency:} EBF acknowledges limitations (e.g., no framework-level RCT) while demonstrating component-level evidence.
\end{enumerate}
\end{tcolorbox}

%%%%%%%%%%%%%%%%%%%%%%%%%%%%%%%%%%%%%%%%%%%%%%%%%%%%%%%%%%%%%%%%%%%%%%%%%%%%%%%
\section{Glossary}
\label{sec:da2-glossary}
%%%%%%%%%%%%%%%%%%%%%%%%%%%%%%%%%%%%%%%%%%%%%%%%%%%%%%%%%%%%%%%%%%%%%%%%%%%%%%%

Key terms used in this appendix. For the comprehensive master glossary, see \textbf{Appendix G (Glossary)}.

\begin{description}[nosep]
    \item[$\mathcal{W}$ (Wissenschaft)] Scientific audience archetype (Ernst Fehr): PhD Economics, formal language, convinced by logic and evidence.

    \item[$\mathcal{P}$ (Praxis)] Practitioner audience archetype (Middle Manager): MBA/experience, concrete examples, convinced by ROI and success stories.

    \item[K/M/L] Response length variants: Kurz (1-2 sentences), Mittel (1 paragraph), Lang (full explanation).

    \item[Epistemische Tags] Parameter provenance markers: EMP (empirical), THR (theoretical), LLM (LLMMC prior), ILL (illustrative), HYP (hypothetical).

    \item[Tier 1/2/3] Evidence hierarchy: Tier 1 = component RCTs, Tier 2 = case documentation, Tier 3 = framework-level evidence.

    \item[$\gamma$] Complementarity parameter measuring interaction effects between utility dimensions.

    \item[$\lambda$] Loss aversion parameter ($\approx 2.25$); losses weighted $\lambda$ times more than equivalent gains.
\end{description}

%%%%%%%%%%%%%%%%%%%%%%%%%%%%%%%%%%%%%%%%%%%%%%%%%%%%%%%%%%%%%%%%%%%%%%%%%%%%%%%
\section{Open Issues}
\label{sec:da2-open}
%%%%%%%%%%%%%%%%%%%%%%%%%%%%%%%%%%%%%%%%%%%%%%%%%%%%%%%%%%%%%%%%%%%%%%%%%%%%%%%

\begin{enumerate}
    \item \textbf{Remaining Questions:} This appendix provides full variants for 4 questions (W1, P1, P6, W3). Full variants for P2-P5, P7-P10, W2, W4-W10 are documented in quick reference form only.

    \item \textbf{Language Consistency:} Response variants mix German and English. Future versions should provide complete translations.

    \item \textbf{Empirical Validation:} The response effectiveness has not been empirically tested. A/B testing of different variants in real presentation contexts is planned.

    \item \textbf{Citation Completeness:} Some LANG versions reference papers not yet in bcm\_master.bib. Integration pending via PIP workflow.
\end{enumerate}

%%%%%%%%%%%%%%%%%%%%%%%%%%%%%%%%%%%%%%%%%%%%%%%%%%%%%%%%%%%%%%%%%%%%%%%%%%%%%%%
\section{References}
\label{sec:da2-references}
%%%%%%%%%%%%%%%%%%%%%%%%%%%%%%%%%%%%%%%%%%%%%%%%%%%%%%%%%%%%%%%%%%%%%%%%%%%%%%%

This appendix draws on the EBF master bibliography containing 2,328 peer-reviewed papers.
Key references for framework defense are integrated throughout the LANG response variants.

\nocite{bcm_master}

\textbf{Core References:}
\begin{itemize}[nosep]
    \item Prospect Theory: Kahneman \& Tversky (1979), Kahneman (1991)
    \item Social Preferences: Fehr \& Schmidt (1999), Fehr \& Fischbacher (2002)
    \item Defaults: Johnson \& Goldstein (2003), Choi et al. (2004), Benartzi \& Thaler (2007)
    \item Social Norms: Cialdini (2006), Allcott (2011), Schultz et al. (2007)
    \item Methodology: Camerer et al. (2003, 2018), Thaler (2015)
    \item Falsificationism: Popper (1959), Kuhn (1962)
\end{itemize}

\end{document}
